Static analysis is a cornerstone of software security assurance for systems code.
Tools such as CodeQL~\cite{codeqlDocs} and Clang Static Analyzer (CSA)~\cite{clangsaDocs} detect vulnerabilities by checking source code against predefined rules without executing it~\cite{engler2001bugs,bessey2010coverity,calcagno2015infer,sadowski2015tricorder}.
The effectiveness of these tools depends on the quality of their vulnerability checkers, each of which must encode the semantic conditions that distinguish a real bug from safe code.
Writing such checkers by hand requires expertise in both the target codebase and the analysis framework, making manual authoring expensive and difficult to scale.

LLM-based patch-driven synthesis offers a promising alternative.
In this paradigm, a security patch serves as a compact before/after specification: the LLM consumes the patch and synthesizes a checker that should fire on the vulnerable revision but remain silent on the fixed revision.
KNighter~\cite{yang2025knighter}, a representative system, has demonstrated that this approach can produce compilable and executable CSA checkers from historical Linux kernel fixes.

However, the resulting checkers may be semantically inadequate.
A security patch is inherently underspecified: a one-line \texttt{goto}-target change may encode an ownership barrier, a widened cast may encode a range guard, and a moved \texttt{free} may encode a lifetime invariant.
A checker that captures only the syntactic shape of the edit can fire on both revisions or miss the vulnerable revision altogether.
Refinement can also remove a useful trigger while suppressing unwanted reports, a failure mode we call over-pruning.

Current refinement strategies, including KNighter's own iterative loop, use compilation and report feedback to guide further edits~\cite{yang2025knighter,madaan2023selfrefine,shinn2023reflexion}.
The checker, patch and reported path help the model reason about the failure.
Even with this context, however, a symptom such as “the vulnerable side is silent” may leave the implementation problem unclear.
A cleanup attribute may name a lowered callback that the checker never recognizes; a cast may hide the expression expected by a callback; or a state lookup may use a different region from the preceding update.

We observe that the analysis framework already exposes information useful for understanding these mismatches.
CSA provides compiler representations and observations of symbolic values, regions and checker callbacks.
CodeQL provides a different interface through its program database and queries~\cite{avgustinov2016ql,codeqlDocs}.
Exploiting this information requires selecting the relevant part of an analyzer's output, relating it to the patch and checker, and distinguishing it from ordinary reports or source text.

We present \toolname, an evidence-guided refinement framework for LLM-generated vulnerability checkers.
\toolname is designed as a plug-in post-generation layer, treating a generated checker as an artifact to be refined rather than regenerated.
First, rule-guided preflight selects patch-relevant evidence requirements.
Second, the collected facts, source context and validation feedback are normalized into a typed evidence bundle, with the origin and unavailable evidence recorded explicitly.
Third, the agent edits the checker and validates it on both revisions.
A retention gate keeps vulnerable-side warning recovery or fixed-side report reduction while rejecting the loss of an existing vulnerable-side warning.

We evaluate the workflow on all 39 selected generated-only CSA checkers from KNighter's artifacts.
Native refinement recovers eight initially silent vulnerable revisions, compared with six without internal evidence.
KNighter's report-driven policy makes no recovery edit on these silent starts.
The native and no-internal variants have nearly equal $F_1$, with 123 and 68 fixed-side reports, respectively.
Among the 15 initially noisy subjects, native refinement instead reduces fixed reports from KNighter's 69 to 44, largely through one subject.
These results led us to examine the actual edits and reports behind the differences.
An audit of all 11 non-tied comparisons, controlled perturbations of four mechanisms and repeated interventions at two checkpoints reveal useful recognition repairs alongside wrong-scope and infeasible-path warnings.
The evidence can help the checker recognize an operation without supplying all the conditions needed to report it correctly.

This paper makes the following contributions:
\begin{itemize}
  \item We present \toolname, a plug-in refinement layer that combines patch facts, source context, compiler representations and checker-execution observations in a typed, provenance-labelled evidence bundle.
  \item We provide a 39-checker comparison with KNighter's refinement loop and a no-internal variant, together with repeated decodes, model sensitivity and a separate CodeQL case.
  \item We analyze the code changes and diagnostic evidence behind the complete non-tie subset, explaining both useful local repairs and the semantic conditions they leave unresolved.
\end{itemize}
